% Figure: temperature-entropy sketch of the ideal Rankine cycle used in the
% plant case study. A dome (saturation curve) with the four state points:
% 1 condenser exit (sat liquid), 2 pump exit, 3 boiler/superheater exit,
% 4 turbine exit (wet vapour). Schematic, not to numerical scale.
\begin{figure}[htbp]
\centering
\begin{tikzpicture}[font=\small, scale=1.0]
  % axes
  \draw[->, thick] (0,0) -- (8.5,0) node[anchor=north east] {specific entropy $s$};
  \draw[->, thick] (0,0) -- (0,6.0) node[anchor=north west, rotate=90, yshift=6mm] {temperature $T$};

  % saturation dome (two-lobe curve via a smooth path)
  \draw[engblue, thick]
    (1.2,0.6) .. controls (1.6,3.2) and (2.6,4.3) .. (3.6,4.55)
    .. controls (4.6,4.35) and (6.4,3.2) .. (7.0,0.6);
  \node[engblue, font=\scriptsize] at (2.0,1.4) {liquid};
  \node[engblue, font=\scriptsize] at (3.6,3.0) {two-phase};
  \node[engblue, font=\scriptsize] at (6.0,1.6) {vapour};

  % state 1: condenser exit, saturated liquid on the dome, low T
  \fill[engblack] (1.45,0.7) circle (1.6pt);
  \node[font=\scriptsize, anchor=east] at (1.40,0.7) {1};
  % state 2: pump exit, slightly above 1 (nearly vertical pump line)
  \fill[engblack] (1.55,1.05) circle (1.6pt);
  \node[font=\scriptsize, anchor=east] at (1.50,1.20) {2};
  % state 3: superheater exit, high T, to the right
  \fill[engblack] (5.6,5.1) circle (1.6pt);
  \node[font=\scriptsize, anchor=south] at (5.6,5.2) {3};
  % state 4: turbine exit, wet vapour on the dome, low T
  \fill[engblack] (6.3,0.85) circle (1.6pt);
  \node[font=\scriptsize, anchor=west] at (6.40,0.85) {4};

  % process lines
  % 1->2 pump (near vertical, isentropic, exaggerated)
  \draw[engred, thick] (1.45,0.7) -- (1.55,1.05);
  % 2->3 boiler + superheat (heat addition, climbing right)
  \draw[engred, thick] (1.55,1.05) .. controls (2.4,1.2) and (3.6,4.0) .. (5.6,5.1);
  \node[engred, font=\scriptsize, anchor=south, rotate=42] at (3.4,2.9) {boiler $q_{\text{in}}$};
  % 3->4 turbine (isentropic expansion, vertical down)
  \draw[englime!70!engblack, thick] (5.6,5.1) -- (6.3,0.85);
  \node[font=\scriptsize, anchor=west, rotate=-78] at (6.15,3.0) {turbine $w_{\text{out}}$};
  % 4->1 condenser (heat rejection, horizontal left at low T)
  \draw[engblue, thick] (6.3,0.85) -- (1.45,0.7);
  \node[engblue, font=\scriptsize, anchor=north] at (3.8,0.72) {condenser $q_{\text{out}}$};
\end{tikzpicture}
\caption[Ideal Rankine cycle on temperature-entropy axes]{The ideal Rankine
cycle of the case study, on temperature-entropy axes. Water leaves the
condenser as saturated liquid at state 1, is pumped to boiler pressure at
state 2 (an almost vertical, low-work step), is heated and superheated to
state 3, expands through the turbine to state 4, and is condensed back to
state 1. The area enclosed is the net work per unit mass; the area under the
$4\to1$ condenser line is the heat the second law forces the plant to reject.
Superheating to state 3 lifts the average temperature of heat addition, which
raises the cycle efficiency toward the Carnot value set by the boiler and
condenser temperatures.}
\label{fig:vol04ch01:rankine-ts}
\end{figure}
